\subsection{Overall Outcomes and Progress (RQ1)}
% All statistics in this section use best-of-three (max per CVE-model pair across 3 runs) -> per-pair average.
Table~\ref{tab:model-results} reports the best-of-three results for each model across 30 CVEs. \texttt{glm-5.2} achieves the highest average overall score (64.6/100), followed by \texttt{mimo-v2.5} (52.9) and \texttt{claude-sonnet-4.6} (50.0). Under the original rubric, 11 of 150 CVE--model pairs achieve near-full P5 credit ($\geq$15/20), while 8 achieve near-full P6 credit. These thresholds describe historical scores, not independently adjudicated end-to-end success.


Historical item annotations identify 16 candidate runs across eight CVE--model groups and five CVEs. The old manuscript describes three execution strategies among these records; target identity, effects, and entry-path scope remain subject to AUD-01. Firstly, same-architecture \textit{chroot} on CVE-2023-26315 (AArch64) enables 6 runs---\texttt{claude-sonnet-4.6} twice, \texttt{deepseek-v4-flash} once, and \texttt{glm-5.2} three times---where the host and target share the same instruction set and no emulation is needed. Secondly, QEMU system mode with MikroTik's CHR image on CVE-2025-6443 enables 6 runs---\texttt{glm-5.2} three times and \texttt{gpt-5.5} three times---where the agent boots the full OS, sends crafted VXLAN packets, and records evidence interpreted as an access-control bypass via packet sniffer evidence. Thirdly, FastCGI socket injection with LD\_PRELOAD stubs on CVE-2023-44418, CVE-2024-5293, and CVE-2020-13389 enables 4 runs by \texttt{glm-5.2}, where the agent builds NVRAM/nofork stub libraries, launches the real \texttt{prog.cgi} or \texttt{httpd} under QEMU user-mode, and triggers SIGSEGV crashes or return-address overwrites on the real binary. In contrast, devices whose firmware is unavailable or requires complex cross-architecture emulation with proprietary NVRAM dependencies almost universally fail at P5.


The leading recorded configuration, \texttt{glm-5.2}, also shows stronger acquisition scores and fallback-plan scores. It averages 11.5/15 on P2 (Firmware Acquisition) under best-of-three scoring, well above \texttt{claude-sonnet-4.6}'s 7.7/15; since P2 gates all downstream phases, this gap cascades---in 5 out of 30 CVE--model pairs, \texttt{claude-sonnet-4.6}'s best plan exceeds 50 yet its best task stalls below 15, all blocked at P2. \texttt{glm-5.2} also achieves the strongest fallback planning (50.8/100 vs.\ 27.5--43.3 for the others), which correlates with deeper pipeline progression. The inspected examples connect continued acquisition attempts and alternative strategies to further progress. These associations do not isolate their causal contribution or an intrinsic property of the model.

Finally, we also report the computational cost of each model in our appendix, including per-run token consumption, API cost, and execution time across all 450 runs.
